% ── CAPÍTULO 2: REQUISITOS DEL SISTEMA ──────────────────────
\section{Requisitos del Sistema}
\label{sec:requisitos}

\subsection{Sistema Operativo}

\begin{table}[H]
    \centering
    \caption{Compatibilidad de sistemas operativos.}
    \begin{tabular}{lll}
        \toprule
        \textbf{Sistema} & \textbf{Versión mínima} & \textbf{Estado} \\
        \midrule
        Windows & 10 / 11 (64-bit) & Verificado \\
        Linux   & Ubuntu 20.04+ o equivalente & Compatible \\
        macOS   & 12 Monterey + & Compatible (sin síntesis de voz) \\
        \bottomrule
    \end{tabular}
\end{table}

\subsection{Python}

Se requiere \textbf{Python 3.10 o superior} (se recomienda 3.12).

\begin{lstlisting}[style=shell, caption={Verificar versión de Python.}]
python --version
# Se espera: Python 3.10.x o superior
\end{lstlisting}

\begin{aviso}
Python 3.9 o inferior no es compatible. Algunas características de tipado estático (\lstinline|X | Y| syntax) requieren Python 3.10+.
\end{aviso}

\subsection{Hardware Recomendado}

\begin{table}[H]
    \centering
    \caption{Recursos de hardware recomendados según el tamaño de la red analizada.}
    \begin{tabular}{llll}
        \toprule
        \textbf{Caso de uso} & \textbf{CPU} & \textbf{RAM} & \textbf{Disco} \\
        \midrule
        Redes pequeñas ($n \leq 10$)  & 2 núcleos & 4 GB  & 500 MB \\
        Redes medianas ($n \leq 15$)  & 4 núcleos & 8 GB  & 1 GB   \\
        Redes grandes ($n \geq 20$)   & 8 núcleos & 16 GB & 5 GB   \\
        \bottomrule
    \end{tabular}
\end{table}

\begin{peligro}
Las redes con $n \geq 20$ y $k = 5$ en modo exhaustivo requieren hasta 16~GB de RAM. Si el sistema no dispone de suficiente memoria, el proceso puede ser terminado por el sistema operativo.
\end{peligro}

\subsection{Dependencias de Python}

Las siguientes bibliotecas se instalan automáticamente al ejecutar la instalación:

\begin{table}[H]
    \centering
    \caption{Dependencias requeridas.}
    \begin{tabular}{lll}
        \toprule
        \textbf{Biblioteca} & \textbf{Versión mínima} & \textbf{Propósito} \\
        \midrule
        numpy        & 2.0.0  & Operaciones matriciales, hipercubos \\
        plotly       & 6.0.0  & Dashboard interactivo HTML \\
        pandas       & 2.0.0  & Lectura y escritura de CSV \\
        colorama     & 0.4.6  & Salida de terminal con colores \\
        pyinstrument & 5.0.0  & Profiling de rendimiento \\
        pyttsx3      & 2.99   & Síntesis de voz (Windows) \\
        \bottomrule
    \end{tabular}
\end{table}

\begin{nota}
El módulo \lstinline|pyttsx3| es opcional en Linux/macOS. Si no está disponible, el software funciona normalmente pero sin el anuncio de voz al finalizar.
\end{nota}

\subsection{Herramienta de Gestión de Paquetes}

Se soportan dos métodos de instalación:

\begin{itemize}[itemsep=4pt]
    \item \textbf{pip}: gestor de paquetes estándar de Python.
    \item \textbf{uv}: gestor moderno y rápido (recomendado). Ver \url{https://docs.astral.sh/uv/}.
\end{itemize}
